% Claim proof file for row claim_3_10 -- "TwoDisjointNode".
% Auto-generated stub by scaffold/solving_current_row.py at row start. A
% manager agent dedicated to writing the TeX proof fills in the body below.
% The proof must:
%   - Stay close to the lecture notes' proof if one exists (always search
%     the LN first for a `\begin{proof}` block immediately following the
%     claim; copy / lightly edit when present).
%   - Otherwise use the LN paradigm and cite prior claims by their `ref`.
%   - Be self-contained enough that a mathematician can follow it without
%     re-reading the LN.
%   - Have no `sorry`, `omitted`, or "obvious" shortcuts.
%
% Compiles standalone via the `subfiles` package.

\documentclass[main]{subfiles}

\begin{document}
% ref: claim_3_10 (proof, with statement restated for readability)
% title: TwoDisjointNode
\def\rowref{claim\_3\_10}
\phantomsection\label{claim_3_10}
%
% Cross-references: see claim_<ref>_statement_<title>.tex in this folder
% for the corresponding statement.

% --- statement (restated) ---------------------------------------------------
\begin{Lem}[Two disjoint node-splitting hard interventions commute]
   Let $G=(J,V,E,L)$ be a CADMG and $W_1, W_2 \ins V$ two disjoint subsets of the output nodes from $G$. 
   Then the CADMG obtained from first node-splitting on $W_1$ and then node-splitting on $W_2$ is the same CADMG that arises from first node-splitting on $W_2$ and then node-splitting on $W_1$:
   \[ \lp G_{\swig(W_1)} \rp_{\swig(W_2)} =  \lp G_{\swig(W_2)} \rp_{\swig(W_1)}  =  G_{\swig(W_1 \dcup W_2)}.   \]
\end{Lem}

% --- proof ------------------------------------------------------------------
% Proof body adapted (light editorial expansion) from the LN at
% lecture-notes/lecture_notes/graphs.tex lines 633 -- 661 -- the
% `\Claude{\begin{proof}...\end{proof}}' block following the claim. The
% outer `\Claude{...}' LN-author wrapper has been stripped; the case
% analysis on the directed-edges paragraph and the bidirected-edges
% paragraph have been written out in full (the LN punts on the latter via
% ``by the same case analysis...''). Cross-refs to \refrow{def_3_12} (for
% the definition of $\swig$) and \refrow{claim_3_7} (the $\spl$-only
% structural mirror) are inserted in prose, never via raw \ref.
\begin{proof}
We show $\lp G_{\swig(W_1)} \rp_{\swig(W_2)} = G_{\swig(W_1 \dcup W_2)}$; the other equality follows by symmetry (swap the roles of $W_1$ and $W_2$). Recall from \refrow{def_3_12} that $G_{\swig(W)}$ is obtained from $G$ by first node-splitting on $W$ and then hard-intervening on the freshly introduced input copies $W^i$; in particular the hard-intervention layer of $\swig$ only deletes directed edges (those with head in $W^i$) and never touches the bidirected-edge set $L$ beyond the relabeling already performed by the node-splitting step.

Since $W_1 \cap W_2 = \emptyset$, each $w \in W_2$ satisfies $w \in V \sm W_1 \ins V_{\swig(W_1)}$, so the inner $\swig(W_2)$ on the left-hand side is well-defined. We write $w^{o_k}, w^{i_k}$ for the copies produced by $\swig(W_k)$ ($k = 1, 2$); for $v \notin W_k$ no copy is created and $v$ is left unchanged. We compare the two CADMGs $\lp G_{\swig(W_1)} \rp_{\swig(W_2)}$ and $G_{\swig(W_1 \dcup W_2)}$ field by field, in the order $(J, V, E, L)$.

\medskip\noindent\emph{Node sets.}
For $G_{\swig(W_1 \dcup W_2)}$ the input and output sets are
\[
   J_{\text{merged}} \;=\; J \dcup W_1^i \dcup W_2^i,
   \qquad
   V_{\text{merged}} \;=\; (V \sm W_1 \sm W_2) \dcup W_1^o \dcup W_2^o.
\]
For $\lp G_{\swig(W_1)} \rp_{\swig(W_2)}$, the inner $\swig(W_1)$ first produces input $J \dcup W_1^i$ and output $(V \sm W_1) \dcup W_1^o$. The outer $\swig(W_2)$ then splits $W_2$, which by disjointness lies entirely in $V \sm W_1 \ins V_{\swig(W_1)}$ and is disjoint from the freshly introduced copies $W_1^o$ (these are not in $V$ to begin with). Hence it adjoins $W_2^i$ to the input set and exchanges $W_2$ for $W_2^o$ inside the output set:
\[
   J_{\text{iter}} \;=\; J \dcup W_1^i \dcup W_2^i,
   \qquad
   V_{\text{iter}} \;=\; ((V \sm W_1) \sm W_2) \dcup W_1^o \dcup W_2^o
   \;=\; (V \sm W_1 \sm W_2) \dcup W_1^o \dcup W_2^o.
\]
Thus $J_{\text{iter}} = J_{\text{merged}}$ and $V_{\text{iter}} = V_{\text{merged}}$.

\medskip\noindent\emph{Directed edges.}
By definition,
\[
   E_{\swig(W_1 \dcup W_2)} \;=\; \lC v_1^i \tuh v_2^o \mid v_1 \tuh v_2 \in E \rC,
\]
where, for $u \in J \cup V$, we set $u^i := u^{i_1}$ if $u \in W_1$, $u^i := u^{i_2}$ if $u \in W_2$, and $u^i := u$ otherwise; the convention for $u^o$ is the analogous one with $o_k$ in place of $i_k$. Disjointness $W_1 \cap W_2 = \emptyset$ makes these three cases mutually exclusive, so the notation is unambiguous.

On the other side,
\[
   E_{\swig(W_1)} \;=\; \lC v_1^{i_1} \tuh v_2^{o_1} \mid v_1 \tuh v_2 \in E \rC,
\]
and applying the outer $\swig(W_2)$ relabels any endpoint lying in $W_2$ to its $i_2 / o_2$ copy. Hence, for each $v_1 \tuh v_2 \in E$, we must check that $(v_1^{i_1})^{i_2} = v_1^i$ and $(v_2^{o_1})^{o_2} = v_2^o$. We do this by case analysis on whether each endpoint lies in $W_1$, in $W_2$, or in neither.

\smallskip
\noindent\emph{Inner endpoint $v_1$ ($i$-superscript).}
\begin{itemize}
  \item $v_1 \in W_1$: then $v_1^{i_1} \in W_1^i \ins J_{\swig(W_1)}$ is a freshly introduced input copy. The outer $\swig(W_2)$ only relabels nodes of $W_2 \ins V$, and $W_1^i$ is by construction disjoint from $V$, so $(v_1^{i_1})^{i_2} = v_1^{i_1}$. This equals $v_1^i$ by the merged convention (since $v_1 \in W_1$).
  \item $v_1 \in W_2$: by disjointness $v_1 \notin W_1$, so $v_1^{i_1} = v_1$ (no relabeling in the inner step). Since $v_1 \in W_2 \ins V \sm W_1 \ins V_{\swig(W_1)}$, the outer $\swig(W_2)$ does relabel: $(v_1)^{i_2} = v_1^{i_2}$, which equals $v_1^i$ (since $v_1 \in W_2$).
  \item $v_1 \notin W_1 \dcup W_2$: $v_1 \notin W_1$ gives $v_1^{i_1} = v_1$, and $v_1 \notin W_2$ gives $(v_1)^{i_2} = v_1$. So $(v_1^{i_1})^{i_2} = v_1 = v_1^i$.
\end{itemize}

\smallskip
\noindent\emph{Outer endpoint $v_2$ ($o$-superscript).}
The same three cases, with $o$ in place of $i$:
\begin{itemize}
  \item $v_2 \in W_1$: $v_2^{o_1} \in W_1^o$ is a fresh output copy, disjoint from $W_2 \ins V$, so the outer $\swig(W_2)$ leaves it alone: $(v_2^{o_1})^{o_2} = v_2^{o_1} = v_2^o$.
  \item $v_2 \in W_2$: disjointness gives $v_2 \notin W_1$ and hence $v_2^{o_1} = v_2$; then $(v_2)^{o_2} = v_2^{o_2} = v_2^o$.
  \item $v_2 \notin W_1 \dcup W_2$: $v_2^{o_1} = v_2$ and $(v_2)^{o_2} = v_2 = v_2^o$.
\end{itemize}

Combining the two endpoint analyses, the assignment $(v_1 \tuh v_2) \mapsto (v_1^{i_1})^{i_2} \tuh (v_2^{o_1})^{o_2}$ coincides with $(v_1 \tuh v_2) \mapsto v_1^i \tuh v_2^o$ for every directed edge of $G$. Therefore
\[
   E_{\lp G_{\swig(W_1)} \rp_{\swig(W_2)}}
   \;=\; \lC v_1^i \tuh v_2^o \mid v_1 \tuh v_2 \in E \rC
   \;=\; E_{\swig(W_1 \dcup W_2)}.
\]
The disjointness hypothesis is invoked exactly in the implications ``$v \in W_2 \Rightarrow v \notin W_1$'' (and symmetrically ``$v \in W_1 \Rightarrow v \notin W_2$''): these ensure that the inner and outer relabeling steps never both fire on the same endpoint, so the iterated superscript is well-defined and matches the merged one.

\medskip\noindent\emph{Bidirected edges.}
Bidirected edges of any CADMG have both endpoints in the output set $V$, and the hard-intervention layer of $\swig$ only modifies the directed-edge set $E$ (it does not touch $L$). Hence, for any $W \ins V$,
\[
   L_{\swig(W)} \;=\; \lC v_1^o \huh v_2^o \mid v_1 \huh v_2 \in L \rC,
\]
i.e.\ the node-splitting step relabels both endpoints of every bidirected edge to their $W^o$ copies, and no $i$-superscript ever appears on a bidirected endpoint. Applying this to $W = W_1 \dcup W_2$ (using the merged $o$-convention) gives
\[
   L_{\swig(W_1 \dcup W_2)} \;=\; \lC v_1^o \huh v_2^o \mid v_1 \huh v_2 \in L \rC,
\]
and applying it twice for the iterated side gives
\[
   L_{\swig(W_1)} \;=\; \lC v_1^{o_1} \huh v_2^{o_1} \mid v_1 \huh v_2 \in L \rC,
\]
after which the outer $\swig(W_2)$ further relabels endpoints in $W_2$ to their $o_2$ copies. The case analysis is identical to the one for the outer endpoint of the directed-edges paragraph, applied separately to \emph{each} endpoint $v_1, v_2$ of every bidirected edge:
\begin{itemize}
  \item $v_k \in W_1$ ($k = 1, 2$): $v_k^{o_1} \in W_1^o$ is a fresh output copy, disjoint from $W_2 \ins V$, hence untouched by the outer $\swig(W_2)$, so $(v_k^{o_1})^{o_2} = v_k^{o_1} = v_k^o$.
  \item $v_k \in W_2$: disjointness gives $v_k^{o_1} = v_k$; then $(v_k)^{o_2} = v_k^{o_2} = v_k^o$.
  \item $v_k \notin W_1 \dcup W_2$: $(v_k^{o_1})^{o_2} = v_k = v_k^o$.
\end{itemize}
Hence
\[
   L_{\lp G_{\swig(W_1)} \rp_{\swig(W_2)}}
   \;=\; \lC v_1^o \huh v_2^o \mid v_1 \huh v_2 \in L \rC
   \;=\; L_{\swig(W_1 \dcup W_2)}.
\]
The bidirected-edge case is structurally simpler than the directed-edge case for two reasons: (i) only the $o$-superscript appears (there is no $i$-superscript on bidirected endpoints), and (ii) the hard-intervention layer of $\swig$ acts trivially on $L$, so the only mutation is the relabeling of output nodes in $W$ to their fresh copies in $W^o$.

\medskip
All four data fields $(J, V, E, L)$ of $\lp G_{\swig(W_1)} \rp_{\swig(W_2)}$ and $G_{\swig(W_1 \dcup W_2)}$ agree, so the two CADMGs are equal. The remaining equality $\lp G_{\swig(W_2)} \rp_{\swig(W_1)} = G_{\swig(W_1 \dcup W_2)}$ follows by the same argument with the roles of $W_1, W_2$ swapped (using $W_1 \dcup W_2 = W_2 \dcup W_1$). This establishes both equalities in the statement.

This is the structural mirror of \refrow{claim_3_7} (the analogous statement for two disjoint $\spl$-only node-splittings): the case analysis is identical modulo the hard-intervention layer of $\swig$, which kills the split edges $v^i \to v^o$ in $E$ but leaves the relabeling on $J$, $V$, the remaining directed edges and $L$ unchanged.
\end{proof}

\end{document}
